\documentclass[fleqn]{article}

\usepackage{a4}
\usepackage{latexsym}
\usepackage{pncproofmacros}
\usepackage{lineno}

\title{An example showing the use of pncproofmacros ...}
\author{W.P.}


\begin{document}

\maketitle

% turning on line-numbers to faciliate review
\linenumbers

%
% Note: the empty lines below should be there explicitly. If you don't 
% like it you'll have to explicitly put a latex new-line command.
%

\begin{PROOF}{main}  

\ASSUME{A1}{a[n]=0}

\ASSUME{A2}{(\forall i: 0\leq i<n: 0<a[j])}

\ASSUME{A3}{0\leq j \leq n}

\GOAL{G1}{0\leq a[j]}

\begin{BODY}

\item[1] \HINT{integer arithmetics}
	\SPACE
	\EXPR{0\leq j \leq n \ \ = \ \ 0\leq j < n \ \vee \ (j=n)}

\item[2] \HINT{rewrite{A2} with 1}
	\SPACE
	\EXPR{0\leq j < n \ \vee \ (j=n)}

\item[3] \HINT{see the subproof below}
	\SPACE
	\EXPR{0\leq j < n \Rightarrow 0\leq a[j]}

  \begin{quote}
  \begin{PROOF}{}

	\ASSUME{A1}{0\leq j < n}

	\ASSUME{G1}{0\leq a[j]}

	\begin{BODY}

	\item[1] \HINT{$\forall$-Elimination on top.A2 using A1}
		\SPACE
		\EXPR{0<a[j]}

	\item[2] \HINT{arithmetics} 
		\SPACE
		\EXPR{0<a[j] \Rightarrow 0\leq a[j]}

	\item[3] \HINT{Modus-ponens on 1,2}
		\SPACE
		\EXPR{0\leq a[j]}

	\end{BODY}

  \end{PROOF}
  \end{quote}

\item[4] \HINT{follows from A1, prove this yourself}
	\SPACE
	\EXPR{(j=n) \Rightarrow 0\leq a[j]}

\item[5] \HINT{Case Split on 2,3,4}
	\SPACE
	\EXPR{0\leq a[j]}

\end{BODY}
\end{PROOF}

\end{document}



